\begin{frontmatter}

\title{Screening for Bid Rigging with Frequent Losers in Public Procurement\texorpdfstring{\footnote{This version: April 2026.}}{}}

\author[insper]{Darcio Genicolo-Martins\corref{cor1}}
\ead{darciogm1@insper.edu.br}
\author[insper]{Paulo Furquim de Azevedo}
\affiliation[insper]{organization={INSPER},
  city={S\~ao Paulo}, country={Brazil}}
\cortext[cor1]{Corresponding author.}

\begin{abstract}
We propose a screening tool for bid-rigging cartels that requires
only participation records and win/loss outcomes---no bid microdata,
no structural assumptions, no enforcement records. The tool identifies
\emph{frequent losers} (FL): firms that never win yet participate in
abnormally many tenders, a pattern consistent with synthetic
competition. Applied to 4.5 million tender-items from Brazilian public
procurement (2009--2019), the screen achieves AUC~$= 0.94$ against
competition-authority convictions, with the data-driven optimal
threshold nearly identical to the baseline rule. FL-flagged
environments show 3.6--7.7\% higher prices in within-item
comparisons and 3.5$\times$ the baseline rate of co-participation
with convicted cartelists. The price gap is a conditional association,
not a causal estimate; supporting diagnostics are coherent with cover
bidding but do not identify it. A flag is a starting point for
investigation, not a verdict.
\end{abstract}

\end{frontmatter}

\noindent\textbf{Keywords:} bid rigging, cover bidding, public procurement, cartel screening, frequent losers

\noindent\textbf{JEL No.:} D44, D73, H57, K21, L41

\bigskip
